% Part: incompleteness
% Chapter: arithmetization-syntax
% Section: proofs-in-ax

\documentclass[../../../include/open-logic-section]{subfiles}

\begin{document}

\olfileid[id]{inc}{art}{pax}
\olsection{\usetoken{P}{derivation} Aksiomatik}

\begin{explain}
Untuk mengaritmetisasi !!{derivation}s aksiomatik, kita harus merepresentasikan
!!{derivation}s sebagai bilangan. Karena !!{derivation}s hanyalah barisan
!!{formula}s, pendekatan yang jelas ialah mengodekan setiap !!{derivation}
sebagai kode barisan kode !!{formula}s di dalamnya.
\end{explain}

\begin{defn}
Jika $\delta$ adalah !!{derivation} aksiomatik yang terdiri atas !!{formula}s
$!A_1$, \dots,~$!A_n$, maka $\Gn{\delta}$ ialah
\[
\tuple{\Gn{!A_1}, \dots, \Gn{!A_n}}.
\]
\end{defn}

\begin{ex}
  Tinjaulah !!{derivation} yang sangat sederhana berikut:
  \begin{derivation}
  1. & $!B \lif (!B \lor !A)$ \\
  2. & $(!B \lif (!B \lor !A)) \lif (!A  \lif (!B \lif (!B \lor !A)))$\\
  3. & $!A  \lif (!B \lif (!B \lor !A))$
  \end{derivation}
  Bilangan G\"odel !!{derivation} ini ialah
  \begin{align*}
  \openTuple\,
    & \Gn{!B \lif (!B \lor !A)}, \\
    &\Gn{(!B \lif (!B \lor !A)) \lif (!A  \lif (!B \lif (!B \lor !A)))},\\
    & \Gn{!A  \lif (!B \lif (!B \lor !A))} \,\closeTuple.
  \end{align*}
\end{ex}

\begin{explain}
Setelah menetapkan suatu representasi bagi !!{derivation}s, kita juga harus
menunjukkan bahwa kita dapat memanipulasi !!{derivation}s tersebut secara
rekursif primitif serta menyatakan sifat dan relasi pentingnya dengan cara yang
sama. Beberapa operasi sederhana: misalnya, jika diberikan bilangan
G\"odel~$d$ dari !!a{derivation}, $(d)_{\len{d}-1}$ memberikan bilangan
G\"odel !!{formula} akhirnya. Sebagian operasi lain jauh lebih sulit.
Setidaknya kita akan membuat sketsa cara menanganinya. Tujuannya ialah
menunjukkan bahwa relasi ``$\delta$ adalah !!a{derivation} dari~$!A$
berdasarkan~$\Gamma$'' bersifat rekursif primitif dalam bilangan G\"odel
$\delta$ dan~$!A$.
\end{explain}

\begin{prop}
\ollabel{prop:followsby}
Relasi-relasi berikut rekursif primitif:
\begin{enumerate}
\item $!A$ adalah aksioma.
\item Baris ke-$i$ dalam $\delta$ dibenarkan oleh modus ponens.
\item Baris ke-$i$ dalam $\delta$ dibenarkan oleh \QR.
\item $\delta$ adalah !!{derivation} yang benar.
\end{enumerate}
\end{prop}

\begin{proof}
Kita harus menunjukkan bahwa relasi-relasi yang bersesuaian antara bilangan
G\"odel !!{formula}s dan bilangan G\"odel !!{derivation}s bersifat rekursif
primitif.
\begin{enumerate}
\item Kita memiliki daftar skema aksioma tertentu, dan $!A$ adalah aksioma
  jika bentuknya diberikan oleh salah satu skema tersebut. Karena daftar skema
  itu berhingga, cukup ditunjukkan bahwa bagi setiap skema aksioma kita dapat
  menguji secara rekursif primitif apakah $!A$ memiliki bentuk tersebut.
  Sebagai contoh, tinjaulah skema aksioma
  \[
  !B \lif (!C \lif !B).
  \]
  $!A$ merupakan instans skema aksioma ini jika terdapat !!{formula}s $!B$
  dan $!C$ sedemikian rupa sehingga kita memperoleh $!A$ dengan merangkaikan
  `$($', lalu $!B$, lalu `$\lif$', lalu `$($', lalu $!C$, lalu `$\lif$', lalu
  $!B$, dan akhirnya~`$))$'. Kita dapat menguji sifat yang bersesuaian dari
  bilangan G\"odel $n$ bagi~$!A$, sebab perangkaian barisan bersifat rekursif
  primitif dan bilangan G\"odel $!B$ serta $!C$ harus lebih kecil daripada
  bilangan G\"odel~$!A$: apabila relasi tersebut berlaku, $!B$ dan $!C$
  keduanya merupakan sub-!!{formula}s dari~$!A$. Jadi, kita dapat
  mendefinisikan:
  \begin{multline*}
  \fn{IsAx}_{!B \lif (!C \lif !B)}(n) \defiff \bexists{b< n}{
    \bexists{c < n}{(\fn{Sent}(b) \land \fn{Sent}(c) \land {}}}\\ n =
  \Gn{(} \concat b \concat \Gn{\lif} \concat \Gn{(} \concat c \concat
  \Gn{\lif} \concat b \concat \Gn{))}).
  \end{multline*}
  Jika kita memiliki definisi semacam itu bagi setiap skema aksioma,
  disjungsi semua definisi tersebut mendefinisikan sifat $\fn{IsAx}(n)$,
  yakni ``$n$ adalah bilangan G\"odel suatu aksioma.''
\item Baris ke-$i$ dalam $\delta$ dibenarkan oleh modus ponens jika dan hanya
  jika terdapat baris $j$ dan $k < i$ dengan !!{sentence} pada baris~$j$
  merupakan suatu formula $!A$, kalimat pada baris~$k$ adalah
  $!A \lif !B$, dan kalimat pada baris~$i$ adalah~$!B$.
  \begin{multline*}
    \fn{MP}(d, i) \defiff \bexists{j < i}{\bexists{k < i}{}}\\
    (d)_k = \Gn{(} \concat (d)_j
      \concat \Gn{\lif} \concat (d)_i \concat \Gn{)}
  \end{multline*}
  Karena kuantifikasi terbatas, perangkaian, dan $=$ bersifat rekursif
  primitif, rumusan ini mendefinisikan relasi rekursif primitif.
\item Suatu baris dalam $\delta$ dibenarkan oleh \QR{} jika berbentuk
  $!B \lif \lforall[x][!A(x)]$, suatu baris sebelumnya adalah
  $!B \lif !A(c)$ untuk suatu !!{constant}~$c$, dan $c$ tidak muncul dalam
  $!B$. Hal ini berlaku jika dan hanya jika
  \begin{enumerate}
  \item terdapat !!a{sentence}~$!B$ dan
  \item !!a{formula}~$!A(x)$ dengan satu-satunya variabel bebas~$x$ sedemikian
    rupa sehingga
  \item baris $i$ memuat $!B \lif \lforall[x][!A(x)]$,
  \item suatu baris $j < i$ memuat $!B \lif \Subst{!A}{c}{x}$ untuk suatu
    konstanta~$c$,
  \item dan konstanta itu tidak muncul dalam~$!B$.
  \end{enumerate}
  Semua hal ini dapat diuji secara rekursif primitif, karena bilangan G\"odel
  $!B$, $!A(x)$, dan $x$ lebih kecil daripada bilangan G\"odel formula pada
  baris~$i$, sedangkan bilangan yang bersesuaian dengan $c$ lebih kecil daripada
  G\"odel formula pada baris~$j$:
  \begin{multline*}
    \fn{QR}_1(d, i) \defiff \bexists{j<i}{\bexists{b < (d)_i}{\bexists{x <
        (d)_i}{\bexists{a < (d)_i}{\bexists{c < (d)_j}{(}}}}}
    \\ \fn{Var}(x) \land \fn{Const}(c) \land {} \\
    (d)_i = \Gn{(} \concat
    b \concat \Gn{\lif} \concat \Gn{\lforall} \concat x \concat a
    \concat \Gn{)} \land {}\\
    (d)_j = \Gn{(} \concat b \concat
    \Gn{\lif} \concat \fn{Subst}(a,c,x) \concat \Gn{)} \land {}\\
    \fn{Sent}(b) \land \fn{Sent}(\fn{Subst}(a,c,x)) \land {}
    \bforall{k < \len{b}}{(b)_k \neq (c)_0}).
  \end{multline*}
  Di sini kita mengandaikan bahwa $c$ dan $x$ adalah bilangan G\"odel dari
  variabel dan konstanta yang ditinjau sebagai suku (jadi, bukan kode
  simbolnya). Kita menguji bahwa $x$ adalah satu-satunya variabel bebas dalam
  $!A(x)$ dengan menguji apakah $\Subst{!A(x)}{c}{x}$ merupakan
  !!a{sentence}, dan memastikan bahwa $c$ tidak muncul dalam $!B$ dengan
  mensyaratkan agar setiap simbol~$!B$ berbeda dari~$c$.

  Versi lain dari \QR{} kita tinggalkan sebagai latihan.
\item $d$ adalah bilangan G\"odel !!{derivation} yang benar jika dan hanya
  jika setiap baris di dalamnya merupakan aksioma atau dibenarkan oleh modus
  ponens atau~\QR. Jadi:
  \[
  \fn{Deriv}(d) \defiff \bforall{i < \len{d}}{(\fn{IsAx}((d)_i) \lor
    \fn{MP}(d,i) \lor \fn{QR}(d, i))}
  \]
\end{enumerate}
\end{proof}

\begin{prob}
Definisikan relasi-relasi berikut seperti dalam
\olref[inc][art][pax]{prop:followsby}:
\begin{enumerate}
\item $\fn{IsAx}_{!A \lif (!B \lif (!A \land !B))}(n)$;
\item $\fn{IsAx}_{\lforall[x][!A(x)] \lif !A(t)}(n)$;
\item $\fn{QR}_{2}(d, i)$ (bagi versi lain dari \QR).
\end{enumerate}
\end{prob}

\begin{prop}
Andaikan $\Gamma$ adalah himpunan rekursif primitif yang terdiri atas
!!{sentence}s. Maka relasi $\Prf[\Gamma](x, y)$ yang menyatakan ``$x$ adalah
kode !!a{derivation}~$\delta$ dari $!A$ berdasarkan~$\Gamma$, dan $y$ adalah
bilangan G\"odel dari~$!A$'' bersifat rekursif primitif.
\end{prop}

\begin{proof}
Andaikan ``$y \in \Gamma$'' diberikan oleh predikat rekursif primitif
$R_\Gamma(y)$. Kita harus menunjukkan bahwa relasi $\Prf[\Gamma](x, y)$
bersifat rekursif primitif, dengan $\Prf[\Gamma](x, y)$ berlaku jika dan hanya
jika $y$ adalah bilangan G\"odel !!a{sentence}~$!A$, dan $x$ adalah kode
!!a{derivation} dari~$!A$ berdasarkan $\Gamma$.

Menurut proposisi sebelumnya, sifat $\fn{Deriv}(x)$, yang berlaku jika dan hanya
jika $x$ adalah kode !!{derivation}~$\delta$ yang benar, bersifat rekursif
primitif. Namun, definisi tersebut belum memperhitungkan himpunan $\Gamma$
sebagai cara tambahan untuk membenarkan baris-baris dalam !!{derivation}.
Pengujian rekursif primitif kita tentang apakah suatu baris dibenarkan oleh
\QR{} juga belum memperhitungkan syarat bahwa konstanta~$c$ tidak boleh muncul
dalam~$\Gamma$. Kita dapat mengubah definisi agar langsung memperhitungkan
$\Gamma$, tetapi lebih mudah menggunakan $\fn{Deriv}$ dan teorema deduksi.
$\Gamma \Proves !A$ jika dan hanya jika terdapat suatu daftar berhingga
!!{sentence}s $!B_1$, \dots, $!B_n \in \Gamma$ sedemikian rupa sehingga
$\{!B_1, \dots, !B_n\} \Proves !A$. Menurut teorema deduksi, hal ini berlaku
jika $\Proves (!B_1 \lif (!B_2 \lif \cdots (!B_n \lif !A)\cdots))$. Kita
dapat menguji secara rekursif primitif apakah !!a{sentence} dengan bilangan
G\"odel~$z$ memiliki bentuk ini. Jadi, alih-alih meninjau $x$ sebagai bilangan
G\"odel !!a{derivation} dari !!{sentence} dengan bilangan G\"odel~$y$
\emph{berdasarkan $\Gamma$}, kita meninjau $x$ sebagai bilangan G\"odel
!!a{derivation} dari suatu implikasi bersarang berbentuk seperti di atas
berdasarkan~$\emptyset$.

Pertama, jika kita memiliki barisan !!{sentence}s, kita dapat secara rekursif
primitif membentuk implikasi yang menjadikan semua kalimat tersebut sebagai
anteseden dan suatu !!{sentence} yang diberikan sebagai konsekuen:
\begin{align*}
  \fn{hCond}(s, y, 0) & = y \\
  \fn{hCond}(s, y, n+1) & = \Gn{(} \concat (s)_{n} \concat \Gn{\lif}
  \concat \fn{hCond}(s, y, n) \concat \Gn{)}\\
  \fn{Cond}(s, y) & = \fn{hCond}(s, y, \len{s})\\
  \intertext{Jadi, kita dapat mendefinisikan $\Prf[\Gamma](x, y)$ melalui}
  \Prf[\Gamma](x, y) & \defiff \bexists{s < \fn{sequenceBound}(x,x)}{(} \\
  &\qquad (x)_{\len{x}-1} = \fn{Cond}(s,y) \land {} \\
  &\qquad\bforall{i<\len{s}}{R_\Gamma((s)_i)} \land {} \\
  &\qquad\fn{Deriv}(x)).
\end{align*}
Batas bagi~$s$ diperoleh dengan memperhatikan bahwa setiap $(s)_i$ merupakan
bilangan G\"odel suatu sub-!!{formula} dari baris terakhir !!{derivation},
yakni lebih kecil daripada $(x)_{\len{x}-1}$. Banyaknya anteseden
$!B \in \Gamma$, yakni panjang~$s$, lebih kecil daripada panjang baris
terakhir~$x$.
\end{proof}

\end{document}
